\section{Proofs for Section~\ref{sec:background}}\label{apdx:proofs}
\subsection{Proof of Proposition~\ref{prop:g_equiv}}
Denoting $w=(dx_1(v),...,dx_d(v))$, we have:

$$\g(v,v)=w^T\cdot \G(p)\cdot w,$$
and 

$$\g_0(v,v)=w^T\cdot w,$$

where $\G(p):=\left[\g_{i,j}\right]_{i,j}$ is the symmetric positive-definite matrix given by $\g_{i,j}=\g(\partial_{i_{|p}},\partial_{j_{|p}})$.\\

Hence,
$$\lambda(p)\cdot \sqrt{\g_0(v,v)}\leq \sqrt{\g(v,v)}\leq \mu(p)\cdot \sqrt{\g_0(v,v)}$$

\subsection{Proof of Proposition~\ref{prop:dis_equiv}}

Let $p\in\man$ and let $U$ be a coordinate neighborhood of $p$.
Up to shrinking $U$, we can assume that:
$$U=\{q\in\man\, ,\, \dis_0(p,q)\leq\eps\},$$
for a given, small enough $\eps>0$, and where $\dis_0$ is the distance associated to the metric $\g_0$, given by 
$$\dis_0(p,q)=\sqrt{x_1(q)^2+\dots+x_d(q)^2},$$
where $q$ is given in local coordinates by $x(q)=(x_1(q),...,x_d(q))$.

From Proposition~\ref{prop:g_equiv}, there exist continuous functions $\lambda,\mu\colon U\rightarrow(0,\infty)$ such that for every $q\in U$ and $v\in\tangent_q\man$, one has
$$\lambda(q)\cdot \sqrt{\g_0(v,v)}\leq \sqrt{\g(v,v)}\leq \mu(q)\cdot \sqrt{\g_0(v,v)}.$$
Furthermore, $\lambda(q),\mu(q)\rightarrow \lambda(p),\mu(p)$ as $q\rightarrow p$ because $\lambda$ and $\mu$ are continuous.

Now, let $\gamma\colon[0,1]\rightarrow\man$ be a piecewise-$C^{\infty}$ curve such that $\gamma(0)=p$ and $\gamma(1)=q\in U$.

\begin{itemize}
\item If $\gamma$ is a segment for $\dis_0$, then it lies in $U$ (as $\dis_0(p,q)\leq\eps$). We also have:
%
\begin{align*}
\dis_0(p,q)&=\int_0^1\sqrt{\g_0\left(\frac{d\gamma(t)}{dt},\frac{d\gamma(t)}{dt}\right)}\, dt\\
&\geq \frac{1}{\sup_{t\in[0,1]}\ \mu(\gamma(t))}\int_0^1\sqrt{\g\left(\frac{d\gamma(t)}{dt},\frac{d\gamma(t)}{dt}\right)}\, dt\\
&=\frac{\length(\gamma)}{\sup_{t\in[0,1]}\ \mu(\gamma(t))}\\
&\geq \frac{\dis(p,q)}{\sup_{t\in[0,1]}\ \mu(\gamma(t))}.
\end{align*}
%
\item If $\gamma$ lies inside $U$, then:
%
\begin{align*}
\length(\gamma)&=\int_0^1\sqrt{\g\left(\frac{d\gamma(t)}{dt},\frac{d\gamma(t)}{dt}\right)}\, dt\\
&\geq \inf_{t\in[0,1]}\ \lambda(\gamma(t)) \cdot \int_0^1\sqrt{\g_0\left(\frac{d\gamma(t)}{dt},\frac{d\gamma(t)}{dt}\right)}\, dt\\
&\geq \inf_{t\in[0,1]}\ \lambda(\gamma(t))\cdot\dis_0(p,q)
\end{align*}
%
\item If $\gamma$ leaves $U$, then there exists a smallest $t_0$ such that $\gamma(t_0)\notin U$. We have:
%
\begin{align*}
\length(\gamma)&\geq\int_0^{t_0}\sqrt{\g\left(\frac{d\gamma(t)}{dt},\frac{d\gamma(t)}{dt}\right)}\, dt\\
&\geq \inf_{t\in[0,1]}\ \lambda(\gamma(t))\cdot\int_0^{t_0}\sqrt{\g_0\left(\frac{d\gamma(t)}{dt},\frac{d\gamma(t)}{dt}\right)}\, dt\\
&\geq \inf_{t\in[0,1]}\ \lambda(\gamma(t))\cdot\eps\\
&\geq \inf_{t\in[0,1]}\ \lambda(\gamma(t))\cdot\dis_0(p,q)
\end{align*}
%\mathieu{Un dessin illustrant les 3 cas pourrait être sympa.}
Like mentioned above, $\lambda(q),\mu(q)\rightarrow \lambda(p),\mu(p)$ as $q\rightarrow p$, therefore (and by shrinking $U$ again if necessary) we can ensure that $\lambda_0=\inf_{r\in U}\lambda(r)>0$ and that $\mu_0=\sup_{r\in U}\mu(r)<\infty$.

We have proved that for every $q\in U$, one has
$$\lambda_0\cdot\dis_0(p,q)\leq \dis(p,q)\leq\mu_0\cdot\dis_0(p,q),$$
and we finally see that, as $q\rightarrow p$:
$$\lambda_0,\mu_0\rightarrow \lambda(p),\mu(p).$$
%Finally, notice that by definition of $\mu_0$ and $\lambda_0$ and since $\lambda(q),\mu(q)\rightarrow 1$ as $q\rightarrow p$, we know that $$\frac{\dis(p,q)}{\dis_0(p,q)}\rightarrow 1\text{ as }q\rightarrow p.$$
\end{itemize}

\subsection{Proof of Proposition~\ref{prop:geo_vol}}
Let us consider a chart $\varphi\colon U\rightarrow\R^d$ around $p$ such that $\varphi(p)=(0,...,0)$ and 

$$\g(\partial_{i_{|p}},\partial_{j_{|p}})=
\begin{cases}
1, \text{ if } i=j\\
0, \text{ otherwise}
\end{cases}$$

so as to have $\g_p=\g_0$, where $\g_0$ is the Euclidean metric in local coordinates.

For a small enough $\eps>0$ and by Proposition~\ref{prop:dis_equiv}, we know that 
$$\ball_0(p,\eps/\mu_0)\subseteq\ball(p,\eps)\subseteq \ball_0(p,\eps/\lambda_0)\subseteq U,$$
for $\lambda_0,\mu_0>0$, where $\ball_0(p,\cdot)$ stands for the Euclidean ball using the distance $\dis_0$ in local coordinates.

As such,
$$\vol\left(\ball_0(p,\eps/\mu_0)\right)\leq \vol\left(\ball(p,\eps)\right)\leq \vol\left(\ball_0(p,\eps/\lambda_0)\right).$$
Furthermore, $\lambda_0,\mu_0\rightarrow 1$ as $\eps\rightarrow 0$ because $\g_p=\g_0$.\\
Denote
$$c=\inf\{\deter(\G(q)),q\in\ball_0(p,\eps/\mu_0)\},$$
$$C=\sup\{\deter(\G(q)),q\in\ball_0(p,\eps/\lambda_0)\},$$
where $\G=\left[\g(\partial_i,\partial_j)\right]_{i,j}$.
Since $\deter(\G(p))=1$ (as $\g_p=\g_0$) and $\G$ is smooth, we have $c,C\rightarrow 1$ as $\eps\rightarrow 0$.

Now, denoting the Lebesgue measure in $\R^d$ as $l$, we have
\begin{itemize}
\item One one hand: \begin{align*}
\vol\left(\ball_0(p,\eps/\mu_0)\right)&=\int_{\{x\in\R^d,\,\Vert x\Vert\leq \eps/\mu_0\}}\sqrt{\deter(\G)}\circ\varphi^{-1}\,dl\\
&\geq \sqrt{c}\int_{\{x\in\R^d,\,\Vert x\Vert\leq \eps/\mu_0\}}\,dl\\
&\geq \sqrt{c}\cdot\alpha_d\cdot \left(\frac{\eps}{\mu_0}\right)^d
\end{align*}

\item On the other hand:\begin{align*}
\vol\left(\ball_0(p,\eps/\lambda_0)\right)&=\int_{\{x\in\R^d,\,\Vert x\Vert\leq \eps/\lambda_0\}}\sqrt{\deter(\G)}\circ\varphi^{-1}\,dl\\
&\leq \sqrt{C}\int_{\{x\in\R^d,\,\Vert x\Vert\leq \eps/\lambda_0\}}\,dl\\
&\leq \sqrt{C}\cdot \alpha_d\cdot \left(\frac{\eps}{\lambda_0}\right)^d
\end{align*}
\end{itemize}

As such,
$$\vol\left(\ball(p,\eps)\right)\underset{\eps\rightarrow 0}{\sim} \alpha_d\cdot\eps^d$$

\subsection{Proof of Proposition~\ref{prop:grad}}
The coordinate vector fields are a basis of the tangent spaces $\tangent_p\man$ for all $p\in U$. As such, in $U$, the gradient $\nabla f$ can be expressed as a linear combination:
$$\nabla f=\sum_{i=1}^d a_i\cdot\partial_i.$$
Now, we know that for all $i$:
$$\g\left(\nabla f,\partial_i\right)=df(\partial_i)=\frac{\partial f}{\partial x_i},$$
and as such
$$\sum_{j=1}^d a_j\cdot \g\left(\partial_j,\partial_i\right)=\frac{\partial f}{\partial x_i}.$$
This gives
$$(a_1,...,a_d)\cdot\left[\g(\partial_i,\partial_j)\right]_{i,j}=\left(\frac{\partial f}{\partial x_1},...,\frac{\partial f}{\partial x_d}\right),$$
and
$$(a_1,...,a_d)=\left(\frac{\partial f}{\partial x_1},...,\frac{\partial f}{\partial x_d}\right)\cdot\left[\g(\partial_i,\partial_j)\right]_{i,j}^{-1}.$$

Furthermore,
\begin{align*}
g\left(\nabla f,\nabla f\right)&=(a_1,...,a_d)\cdot\left[\g(\partial_i,\partial_j)\right]_{i,j}\cdot(a_1,...,a_d)^T\\
&=\left(\frac{\partial f}{\partial x_1},...,\frac{\partial f}{\partial x_d}\right)\cdot\left[\g(\partial_i,\partial_j)\right]_{i,j}^{-1}\cdot\left[\g(\partial_i,\partial_j)\right]_{i,j}\cdot \left[\g(\partial_i,\partial_j)\right]_{i,j}^{-1}\cdot \left(\frac{\partial f}{\partial x_1},...,\frac{\partial f}{\partial x_d}\right)^T\\
&=\left(\frac{\partial f}{\partial x_1},...,\frac{\partial f}{\partial x_d}\right)\cdot \left[\g(\partial_i,\partial_j)\right]_{i,j}^{-1}\cdot \left(\frac{\partial f}{\partial x_1},...,\frac{\partial f}{\partial x_d}\right)^T.
\end{align*}
As such,
$$\frac{1}{\mu(p)}\cdot \sqrt{\sum_{i=1}^d\frac{\partial f(p)}{\partial x_i}^2}\leq \sqrt{\g\left(\nabla f(p),\nabla f(p)\right)}\leq \frac{1}{\lambda(p)}\cdot \sqrt{\sum_{i=1}^d\frac{\partial f(p)}{\partial x_i}^2}.$$




%\subsection{Proof of Lemma~\ref{lem:F}}
%The proof for this lemma follows closely from the proof of Theorem 3.2. in \cite{milnor1963morse}. We will construct a function $F_{\epsilon}$ that coincides with $f$ everywhere but on a small neighborhood around the critical point $p$.\\
%By Lemma~\ref{lem:morse}, there exists a Morse chart $\varphi\colon U\rightarrow\mathbb{R}^n$ where $f=v-\xi+\eta$, such that $\forall q\in U_i$:
%\begin{itemize}
%    \item $\xi(q)=\sum_{j=1}^{i}\varphi(q)_j^2$,
%    \item $\eta(q)=\sum_{j=i+1}^{n}\varphi(q)_j^2$,
%\end{itemize}
%where $i$ depends only on the Hessian of $f$ at $p$.\\
%
%We also have 
%
%$$\ball_{\sqrt{\epsilon}}=\left\{q\in U,\, \xi(q)+\eta(q)\leq\epsilon\right\}.$$
%
%Let $\mu\colon \mathbb{R}\rightarrow\mathbb{R}$ be a smooth function such that:
%\begin{itemize}
%    \item $\mu(0)>\epsilon$,
%    \item $\mu(r)=0$ if $r\geq \epsilon$, and
%    \item $-1<\mu'(r)\leq 0$.
%\end{itemize}
%Now, define $F_\epsilon$ to be:
%\begin{align*}
%    F_\epsilon\colon\man&\longrightarrow\mathbb{R}\\
%    q&\longmapsto \begin{cases}
%    f(q),& \text{ if } q\notin U\\
%    f(q)-\mu(\xi(q)+\eta(q)),& \text{ if } q\in U
%    \end{cases}
%\end{align*}
%
%$F_\epsilon$ is smooth and $F_\epsilon\leq f$. Moreover, outside $\ball_{\sqrt{\epsilon}}$, $f$ and $F_\epsilon$ coincide.\\
%
%Notice also that the critical points of $F_\epsilon$ are the same as those of $f$. This is because the two functions coincide outside of $U$, and inside $U$ we have:
%$$\nabla F_\epsilon=\left(-1-\mu'(\xi+\eta)\right)\cdot\nabla\xi+\left(1-\mu'(\xi+\eta)\right)\cdot\nabla\eta,$$
%$$\left(-1-\mu'(\xi+\eta)\right)<0,$$
%$$\left(1-\mu'(\xi+\eta)\right)\geq 1,$$
%and $\nabla\xi$ and $\nabla\eta$ vanish only at the origin of the chart which corresponds to $p$.\\\\
%Now, $$F_\epsilon^{-1}[f(p)-\epsilon,f(p)+\epsilon]\subseteq f^{-1}[f(p)-\epsilon,f(p)+\epsilon].$$
%
%This is because $f$ and $F$ coincide outside of $\ball_{\sqrt{\epsilon}}$ and inside we have $f\geq F$ and
%$$f=f(p)-\xi+\eta\leq f(p)+\xi+\eta\leq f(p)+\epsilon.$$
%
%Hence, $F_\epsilon^{-1}[f(p)-\epsilon,f(p)+\epsilon]$ can not contain any critical points aside from $p$. However,
% $$F_\epsilon(p)=f(p)-\mu(0)<f(p)-\epsilon.$$
%Therefore, $F_\epsilon^{-1}[f(p)-\epsilon,f(p)+\epsilon]$ does not contain critical points.
% 
%
